\documentclass[book.tex]{subfiles}

\begin{document}

\part{Differentiable Programming}
\label{part:diffprog}
%\end{center}

In this Part of the book, we cover the foundations of differentiable programming. We start by providing an introduction to Numerical methods. These techniques form the foundation of differentiable programming. We review standard numerical methods with an eye towards constructing a library of differentiable numerical primitives that can be reused and extended to create more sophisticated architectures.

We then move on to covering foundational topics in machine learning and introduce foundational deep architectures such as multilayer perceptrons, convolutional networks, recurrent networks, and transformers among others.

We end the Part by moving into the theory of programming languages. The mathematical question of which programs can be successfully differentiated is still emerging. We introduce the basics of type theory and programming language theory and cover the new but foundational ideas surrounding the differentiable lambda calculus.

All the code in this Part of the book is written in Physika, our new differentiable dependently typed programming language. Physika combines modern advances in automatic differentiation with strong type systems to create a rich programming language equally capable of modeling physical systems and differentiable models.

\subfile{part_prog_ProgrammingLanguages}

%\subfile{part_prog_NumericalMethods}
\subfile{part_prog_DifferentiableModels}

% equivariant nets 


\end{document}